\documentclass[aps,pop,preprint]{revtex4}
\usepackage{amsmath}
\usepackage{amssymb}

\newcommand{\Rm}{\mathrm{Rm}}
\renewcommand{\Re}{\mathrm{Re}}
\newcommand{\xhat}{\hat{\mathbf{x}}}
\newcommand{\yhat}{\hat{\mathbf{y}}}
\newcommand{\zhat}{\hat{\mathbf{z}}}

\begin{document}
	% \title{Model equations}
	% \begin{abstract}
	% 
	% \end{abstract}
	% \maketitle
	
	\section{Model Derivation}
	\label{sec:derivation}
	We start from the MHD equations for an incompressible, unstratified plasma with finite viscosity and resistivity:
	\begin{equation}
	\rho_0 \left( \frac{\partial}{\partial t} + \mathbf{v} \cdot \nabla \right) \mathbf{v} = -\nabla p + \frac{1}{c} \mathbf{J} \times \mathbf{B} + \rho_0 \nu \nabla^2 \mathbf{v}
	\end{equation}
	\begin{equation}
	\frac{\partial}{\partial t} \mathbf{B} = \nabla \times (\mathbf{v} \times \mathbf{B}) + \eta \nabla^2 \mathbf{B}
	\end{equation}
	\begin{equation}
	\nabla \cdot \mathbf{v} = 0
	\end{equation}
	\begin{equation}
	\nabla \cdot \mathbf{B} = 0.
	\end{equation}
	We're assuming we have an initial flow of the form $\mathbf{V} = U_0 \tanh (z/d) \hat{\mathbf{x}}$ and an initial field $\mathbf{B} = B_0 \hat{\mathbf{x}}$. 
	As our non-dimensionalization, use $d$, $U_0$, and $B_0$. 
	Then the first two equations can be written as  % this intermediate step is from memory/written hastily, so I may have left something out, double-check later
	\begin{equation}
	\left( \frac{\partial}{\partial t} + \mathbf{v} \cdot \nabla \right) \mathbf{v} = -\nabla p + \frac{1}{M_A^2} \mathbf{J} \times \mathbf{B} + \frac{1}{\Re} \nabla^2 \mathbf{v}
	\end{equation}
	and
	\begin{equation}
	\frac{\partial}{\partial t} \mathbf{B} = \nabla \times (\mathbf{v} \times \mathbf{B}) + \frac{1}{\Rm} \nabla^2 \mathbf{B},
	\end{equation}
	where now everything is nondimensionalized, $p$ is really $p/\rho_0$, and $M_A \equiv U_0/v_A \propto U_0/B_0$ ($v_A$ is the Alfv\'{e}n velocity). 
	So a weak magnetic field means a high $M_A$, and the threshold for stabilizing KH in ideal MHD ($\nu = 0$, $\eta = 0$) is in the neighborhood of $M_A \approx 1$.
	
	\subsection{2D case} \label{derivation:subsec:2d}
	In the 2D case we may define our streamfunction $\phi$ and flux function $\psi$ such that $\mathbf{v} = \hat{\mathbf{y}} \times \nabla \phi$ and $\mathbf{B} = \hat{\mathbf{y}} \times \nabla \psi$. 
	Then the equations become:
\begin{equation}
\frac{\partial}{\partial t} \nabla^2 \phi + \left\{\nabla^2 \phi, \phi \right\} = \frac{1}{M_A^2}\left\{\nabla^2 \psi, \psi \right\} + \frac{1}{\mathrm{Re}}\nabla^4 \phi
\end{equation}
\begin{equation}
\frac{\partial}{\partial t} \psi = \left\{ \phi, \psi \right\} + \frac{1}{\mathrm{Rm}}\nabla^2 \psi
\end{equation}
(where $\left\{f, g\right\} \equiv \partial_x f \partial_z g - \partial_z g \partial_x f$).
\subsection{3D Case}
For the 3D case, we'll use the vector potential $\mathbf{B} = \nabla \times \mathbf{A}$ in the Coulomb gauge such that $\nabla \cdot \mathbf{A} = 0$ and the induction equation reads
	\begin{equation}
	\frac{\partial  \mathbf{A}}{\partial t} + \nabla \Phi = \mathbf{v} \times \mathbf{B} + \frac{1}{\Rm} \nabla^2 \mathbf{A}.
	\end{equation}
We'll need to expand all the non-linear terms into their quasi-linear forms. We'll use capital letters for background quantities and lower case for fluctuations, $\mathbf{v}_{tot} = \mathbf{U} + \mathbf{u}$,

Noting that the background current density $\mathbf{J} = \nabla \times \mathbf{B} = 0$, the Lorentz force becomes
\begin{equation}
    \mathbf{J_{tot} \times B_{tot}} = \mathbf{j \times B} = j_z B_0 \yhat + j_y B_0 \zhat
\end{equation}
with
\begin{equation}
    \mathbf{j} =  (\partial_y b_z - \partial_z b_y) \xhat + (\partial_z b_x - \partial_x b_z) \yhat + (\partial_x b_y - \partial_y b_x) \zhat.
\end{equation}

Next, we need the non-linear term in the induction equation, 
\begin{equation}
    \mathbf{v_{tot} \times B_{tot}} = \mathbf{U \times b} + \mathbf {u \times B} = v B_0 \xhat + w B_0 \yhat.
\end{equation}
Finally, we need the non-linear advection term,
\begin{equation}
    \mathbf{v_{tot}}\cdot \nabla \mathbf{v_{tot}} = \mathbf{U} \cdot \nabla \mathbf{u} + \mathbf{u} \cdot \nabla \mathbf{U} = U \partial_x \mathbf{u} + w \partial_z U \xhat.
\end{equation}

Breaking these into component form, and returning to non-dimensional form, we arrive at our final equation set for Dedalus,
\begin{align}
    \partial_t u + U\partial_x u + w \partial_z U + \partial_x p - \frac{1}{\Re} \nabla^2 u &= 0\\
    \partial_t v + U\partial_x v + \partial_y p - \frac{j_z}{M_A^2} - \frac{1}{\Re} \nabla^2 v &= 0\\
    \partial_t w + U\partial_x w + \partial_z p - \frac{j_y}{M_A^2} - \frac{1}{\Re} \nabla^2 w &= 0\\
    \partial_x u + \partial_y v + wz &= 0\\
    \partial_t A_x + \partial_x \Phi - v - \frac{1}{\Rm} \nabla^2 A_x &= 0\\
    \partial_t A_y + \partial_y \Phi - w - U b_z - \frac{1}{\Rm} \nabla^2 A_y &= 0\\
    \partial_t A_z + \partial_z \Phi - U b_y - \frac{1}{\Rm} \nabla^2 A_z &= 0\\
    \partial_x A_x + \partial_y A_y + A_zz &= 0,
\end{align}
not including the first order reduction variables $uz, vz,A_xz, A_yz,$
\end{document}